\section{Related Work}
\label{sec:related}

\paragraph{Continual Knowledge Graph Embedding.}
Continual learning for knowledge graphs has received growing attention as real-world KGs evolve over time.
LKGE~\citep{cui2023lifelong} provides a modular framework combining embedding models (TransE, DistMult, RotatE) with continual strategies (EWC, experience replay, knowledge distillation), establishing competitive baselines for continual KGE.
PS-CKGE~\citep{wu2024ps} introduces prototype-guided selective consolidation, identifying parameters critical to past knowledge through entity prototypes.
IncDE~\citep{song2023incde} handles incremental KG embedding by decomposing entity representations into historical and incremental components.
Daruna et al.~\citep{daruna2021continual} study continual graph learning for robotics, demonstrating that topology-aware replay selection outperforms random sampling.
All these methods operate exclusively on graph structure and apply uniform regularization, ignoring the multimodal attributes that modern biomedical KGs provide.

\paragraph{Multimodal Graph Learning.}
Integrating heterogeneous modalities into graph representations has shown promise across domains.
MSCGL~\citep{liang2022foundations} explores multimodal fusion for graph-structured data, though without continual learning considerations.
Recent work on multimodal KG embedding~\citep{liang2022foundations} demonstrates that textual and visual features can complement structural patterns, improving link prediction accuracy.
However, these methods assume a static graph setting and do not address the catastrophic forgetting problem that arises when the graph evolves.
\ours bridges this gap by combining multimodal encoding with modality-aware continual learning.

\paragraph{Biomedical Knowledge Graph Reasoning.}
PrimeKG~\citep{chandak2023building} integrates 20+ biomedical databases into a unified graph with 129K+ nodes and 8.1M+ edges across 10 node types and 30 edge types, providing rich multimodal annotations including textual descriptions and molecular structures.
TxGNN~\citep{huang2024zero} leverages PrimeKG for zero-shot drug repurposing via graph neural networks, demonstrating the value of multi-relational biomedical reasoning.
Both approaches treat the KG as static, retraining from scratch when data changes.
Our work extends biomedical KG reasoning to the continual setting with explicit multimodal awareness.

\paragraph{Elastic Weight Consolidation.}
EWC~\citep{kirkpatrick2017overcoming} prevents catastrophic forgetting by penalizing changes to parameters that are important for previously learned tasks, as measured by the diagonal of the Fisher information matrix.
Zenke et al.~\citep{zenke2017continual} propose Synaptic Intelligence (SI), which accumulates parameter importance online during training.
Li and Hoiem~\citep{li2017learning} introduce Learning without Forgetting (LwF), using knowledge distillation from the previous model as a soft regularizer.
Standard EWC computes a single Fisher matrix over all parameters; our modality-aware extension maintains separate Fisher matrices per encoder module, enabling heterogeneous regularization.

\paragraph{Experience Replay for Continual Learning.}
Replay-based methods store and revisit a subset of past examples to mitigate forgetting.
Rolnick et al.~\citep{rolnick2019experience} show that even a small replay buffer substantially reduces forgetting in deep networks.
GEM~\citep{lopez2017gradient} and A-GEM~\citep{chaudhry2019efficient} use stored gradients as constraints during optimization.
For knowledge graphs, BER-style replay~\citep{cui2023lifelong} stores representative triples from previous tasks.
Our multimodal memory replay extends this by using K-means clustering on structural embeddings to ensure topological diversity in the selected exemplars.

\paragraph{LLMs for Knowledge Graphs.}
Pan et al.~\citep{pan2024unifying} survey the intersection of large language models and knowledge graphs, identifying opportunities for LLMs to enhance KG completion, question answering, and reasoning.
Retrieval-augmented generation (RAG)~\citep{lewis2020retrieval} enables LLMs to ground responses in retrieved KG triples, though at the cost of inference-time retrieval latency.
While LLM-based approaches show promise for KG question answering, they struggle with the structured link prediction tasks that require scoring all candidate entities---a setting where embedding-based methods remain dominant.
